\documentclass[11pt]{article}

% Package Header

\usepackage[paperheight=11in,
			paperwidth=8.5in,
			inner=0.4in,
			textheight=8.9in,
			textwidth=7.8in,
			marginparwidth=120pt,
			marginparsep=32pt,
			bottom=1.4in,
			footskip=1.0in]{geometry}
\geometry{letterpaper}                   % ... or a4paper or a5paper or ...

\usepackage{array}
\usepackage{amssymb}
\usepackage{asymptote}
\usepackage{amsmath}
\usepackage{amsfonts}
\usepackage{animate}
\usepackage{arrayjobx}
\usepackage{arydshln}
\usepackage{bookman}
\usepackage{color}
%\usepackage{colortbl}
\usepackage{cancel}
\usepackage{calculator}
\usepackage{collcell}
\usepackage{epstopdf}
\usepackage{enumerate}
\usepackage{epsfig}
%\usepackage{etex}
\usepackage[english]{babel}
\usepackage{etoolbox}
% \usepackage{fancyhdr}
\usepackage{fancybox}
\usepackage[nomessages]{fp}
\usepackage{graphics}
\usepackage{graphicx}
\usepackage{ifthen}
\usepackage{ifmtarg}
\usepackage[latin1]{inputenc}
\usepackage{latexsym}
\usepackage{multirow}
\usepackage{multicol}
\usepackage{multimedia}
\RequirePackage[noplaybutton]{media9}
\usepackage{mathbbol}
\usepackage{mathtools}
\usepackage{multido}
\usepackage{palatino}
\usepackage{polynom}
\usepackage{pxfonts}
\usepackage{pifont}
\usepackage{pgffor}
%\usepackage{pgfplots}
\usepackage{pdfpages}
\usepackage[parfill]{parskip}
\usepackage{soul}
\usepackage{tikz}
\usetikzlibrary{arrows,shapes,trees,shapes.geometric,calc}
\usepackage{txfonts}
\usepackage{tcolorbox}
\usepackage{wasysym}
\usepackage{wrapfig}
\usepackage{xcolor}
\usepackage{xmpmulti}
\usepackage{xparse}
\usepackage{xstring}
\usepackage{enumitem}

%\usepackage{hyperref}
\usepackage{geometry}
\usepackage{amsmath}
\usepackage{amssymb}
\usepackage{amsthm}
\usepackage{amscd}
\usepackage{amsfonts}
\usepackage{graphicx}%
\usepackage{fancyhdr}
\usepackage{array}
\usepackage{pgf,tikz}
\usetikzlibrary{arrows}

% Uncomment the following lines to use the Palatino font.  Remove the
% [osf] bit if you don't like the old style figures.

%\usepackage[T1]{fontenc}
\usepackage[osf]{mathpazo}
%\usepackage{mathpazo}

\geometry{textheight=10.0in, textwidth=7.7in}

%% Customize page headers
%\pagestyle{myheadings}
%\markright{\name}
%\thispagestyle{empty}

% Customize section headings
\usepackage{sectsty}
\subsectionfont{\rmfamily\mdseries\itshape\large}

% Don't indent paragraphs.
\setlength\parindent{0em}


\begin{document}

\input{../../../../../Headers/TestDefinitionHeader.tex}

$\quad$

\vspace{1.0cm}

\begin{center}
\LARGE {MA 301} \\

\Large{Test 1} \\

\large{Fall - 2014} \\

\vspace{0.5cm}

\setlength{\extrarowheight}{0.4cm}
\begin{tabular}{|c|c|c|}
\hline
Problem & points & out of  \\
\hline
\hline
1 & & 7 \\
\hline
2 & & 7 \\
\hline
3 & & 7 \\
\hline
4 & & 4 \\
\hline
5 & & 7 \\
\hline
6 & & 8 \\
\hline
7 & & 2 \\
\hline
EC & & 5 \\
\hline
\hline
Score & & 42\\
\hline
\end{tabular}
\vspace{0.5cm}

%--------TABLE---------
	\setlength{\extrarowheight}{0.1cm}
	\begin{tabular}{rcc}
		Name: & \underline{$\qquad \qquad \qquad \qquad \quad$}& \\
		&&$\hspace{1.4cm}$\\
		HR: & \underline{$\qquad \qquad \qquad \qquad \quad$}& \\
		&&\\
		Signature: & \underline{$\qquad \qquad \qquad \qquad \quad$}& \\
	\end{tabular}
%--------------------

\vspace{0.5cm}

	\setlength{\extrarowheight}{0.1cm}
	\begin{tabular}{| p{12cm} |}
		\hline
		\multicolumn{1}{|c|}{{\bf READ BEFORE YOU START}} \\
		\hline
		 * Show all work.\\
		 * Change your calculator mode to RADIANS.\\
		 * Sloppy solutions may be penalized.\\
		 * Use calculators ONLY for computations.\\
		 * You must write up every problem in a clear and \\
		 $\,\,$ organized way.\\
		 * If you write solutions on a page different than the \\
		 $\,\,$ question, then clearly state where it is located.\\
		\hline
	\end{tabular}

\end{center}

\vspace{1.0cm}

\thispagestyle{empty}
\newpage
\thispagestyle{empty}

{\bf Problem 1:} Let {\bf F}$=[z^2,x-y,7]$ and $C$ be the intersection of the surfaces $x=y+z$ and $y=z^2$ with $0 \le y \le 9$.

\vspace{0.1cm}

(a) \quad Compute \,\,$\D{\int \limits_{C} {\bf F} \cdot d{\bf r}}$

\vspace{10.2cm}

(b) \quad Completely set-up, but DO NOT compute the integral needed to find the length of $C$.  (arclength = $\int \limits_{C} 1 \, ds$)

\thispagestyle{empty}
\newpage
\thispagestyle{empty}

{\bf Problem 2:} Let \, $\D{{\bf F} = \left[{y-\ln(y), -\frac{x}{y}+x, 1}\right]}$ compute the work required to move an object along $C$ (shown below). Justify your work.

\hspace{14.0cm}	\includegraphics[scale = 0.38]{../../../Figures/Test1_Fig3.png}

\thispagestyle{empty}
\newpage
\thispagestyle{empty}

{\bf Problem 3:} Let $\D {\bf F}(x,y) = (xe^{y}-3xy) {\bf i} + \frac{3}{2}y^2 {\bf j}$ and suppose \(C\) consist of the following 3 curves: the line segment from $(0,0)$ to $(2,0)$ then the curve $y=4-x^2$ from $(2,0)$ to $(0,4)$ and  the line segment from $(0,4)$ to $(0,0)$.

\vspace{0.1cm}

(a) \quad Sketch $C$

\vspace{0.01cm}

(b) \quad Compute \(\D{\int_C {\bf F}\cdot d{\bf r}}\)

\thispagestyle{empty}
\newpage
\thispagestyle{empty}

{\bf Problem 4:} Suppose a straight wire from initial point \((0,1,5)\) to end point \((2,2,3)\) has a linear density function given by $\delta(x,y,z) = x^3 z$. Find the total mass of the wire.

%\thispagestyle{empty}
%\newpage
%\thispagestyle{empty}
%
%{\bf Problem 4:} Let $R$ be a region bounded by the upper half of the circle $x^2+y^2=10$ and the x-axis.  Let $C$ be the boundary of $R$ with counter-clockwise orientation.  Compute the following line integral 
%\begin{center}
%	$\D{\int_{C} (15x-4xy) \, dx + (\cos y-6xy) \, dy}$ 
%\end{center}

\thispagestyle{empty}
\newpage
\thispagestyle{empty}

{\bf Problem 5:} Answer the following true/false questions

	\setlength{\extrarowheight}{-0.2cm}
	\begin{tabular}{p{0.2cm}cp{16.2cm}}
		(a) & True \, or \, False: & If $C$ is an ellipse and {\bf F} is irrotational then $\D{\int \limits_{C} {\bf F} \cdot d{\bf r}=0}$\\
		\\
%		True \, or \, False & If $\WORK = 0$ for one particular closed path, the ${\bf F}$ is path-independent.\\
%		\hline
%		True \, or \, False & If ${\bf F}$ is path-independent, then there is a potential function for ${\bf F}$\\
%		\hline
		(b) & True \, or \, False: & If $\WORK \ne 0$ for some closed path $C$, then ${\bf F}$ is conservative.\\
		\\
		(c) & True \, or \, False: & If ${\bf F} = \nabla f$ then ${\bf F}$ is conservative.\\
		\\
%		True \, or \, False & If ${\bf F}$ is incompressible then ${\bf F}$ is conservative.\\
%		\hline
		(d) & True \, or \, False: & If ${\bf F}$ is conservative then the work required to move an object along any path is zero.\\
		\\
		(e) & True \, or \, False: & A vector field, ${\bf F}$, can be both incompressible and irrotational.\\
		\\
		(f) & True \, or \, False: & If ${\bf F}$ is path-independent and $C$ is any closed curve, then $\WORK = 0$\\
		\\
		(g) & True \, or \, False: & Line integrals over all constant vector fields are path-independent.\\
		\\
		\hline
%		True \, or \, False & If $S_1$ is a rectangle with area = 1 and $S_2$ is a rectangle with area = 2, then
%$2\iint \limits_{S_1} {\bf F} \cdot d{\bf A} = \iint \limits_{S_2} {\bf F} \cdot d{\bf A}$\\
%		\hline
%		True \, or \, False & Let $S$ be a sphere of radius 1 centered at the origin and oriented outward.  If ${\bf F} = [x,y,z]$ then $\FLUX$ is positive\\
%		\hline
%		True \, or \, False & Let $S$ be a sphere of radius 1 centered at the origin and oriented outward.  If $\FLUX = 0$, then ${\bf F} = {\bf 0}$\\
%		\hline
%		True \, or \, False & Let $S$ be the surface of a cube centered at the origin with side length 2 and oriented outward.  If ${\bf F} = 2\VI - \VK$ then the flux of ${\bf F}$ through $S$ is 27.\\
%		\hline
%		True \, or \, False & If ${\bf F}$ and ${\bf G}$ are vector-fields satisfying div(${\bf F}$) = div(${\bf G}$), then ${\bf F} = {\bf G}$\\
%		\hline
%		True \, or \, False & If ${\bf F}$ is a vector-field such that div(${\bf F}$) = 1 and $S$ is a closed surface oriented outward, then $\FLUX$ is equal to the volume enclosed by $S$.\\
%		\hline
%		True \, or \, False & If $\FLUX = 12$ and $S$ is a flat disk of area $4\pi$, then div$({\bf F}) = \frac{3}{\pi}$\\
%		\hline
	\end{tabular}

\vspace{0.2cm}

{\bf Problem 6:} Let {\bf F} be the vector field and and $C_1$ and $C_2$ be the curves shown in the provided figure.

%--------TABLE---------
\setlength{\extrarowheight}{-0.2cm}
\begin{tabular}{p{0.4cm}p{10.4cm}c}
	\\
	(a) & Is {\bf F} incompressible? &\multirow{5}{*}{\includegraphics[scale = 0.3]{../../../Figures/Test1_Fig5.png}}\\
	\\
	(b) & Is {\bf F} irrotational? &\\
	\\
	(c) & Is {\bf F} conservative? &\\
	\\
	(d) & Would the work for a mass traveling along $C_1$ be positive, negative, or zero? Explain. & \\
	\\\\\\\\\\\\\\\\
	(e) & Would the work for a mass traveling along $C_2$ be positive, negative, or zero? Explain. &\\
	\\\\\\\\\\\\\\\\
	(f) & Which is larger: \,\, $|\text{curl}({\bf F}(2,2))|$ \, or \, div(${\bf F}(2,2)$)\,? Explain. &\\
	\\\\\\\\\\\\\\\\
	\hline
\end{tabular}

{\bf Problem 7:}  Let $f(x,y,z)$ be a scalar-valued function.  Using the concepts from this class, explain why curl$(\D \nabla f) = {\bf 0}$ no matter what $f(x,y,z)$ is.  If you cannot explain it, then prove it mathematically.

\thispagestyle{empty}
\newpage
\thispagestyle{empty}

{\bf Extra Credit:} Let $C$ be a simple closed curve enclosing a region $R$. 

\vspace{0.2cm}

(a) \quad Use Green's Theorem to show \quad $\D \oint \limits_{C} x dy = Area(R)$ 

\vspace{5.2cm}

(b) \quad Use this to find the area of the region with boundary given by ${\bf r}(t) = [\sin 2t,2\sin t]$ with $0 \le t \le \pi$.

%--------TABLE---------
	\setlength{\extrarowheight}{0.2cm}
	\begin{tabular}{p{0.6cm}p{11.2cm}p{1cm}c}
		{\it Hint:} & {\it The trig identity: $\sin(2t) = 2\sin t \cos t$ \, might come in handy.} & & \multirow{2}{*}{\includegraphics[scale = 0.3]{../../../Figures/Test1_Fig9.png}} \\
		&&&\\
	\end{tabular}
%--------------------

\thispagestyle{empty}
%\newpage
%\thispagestyle{empty}
%
%{\bf Problem 4:}  Let \({\bf F}=\nabla (2x^{2}+2y^{3})\) and  \(C\) is the quarter of the circle \(x^2+y^2 = 1\) in the fourth quadrant, oriented counterclockwise. Compute \(\D{\int_C {\bf F}\cdot d{\bf r} }\)

%{\bf Problem 5:} Answer the following questions.  
%
%(a) Consider the vector field \({\bf F}\) shown in the figure, together with the paths \(C_1\), \(C_2\), and \(C_3\). Order the line integrals $\D{\left(\int_{C_1} {\bf F} \cdot d{\bf r}, \,\, \int_{C_2} {\bf F} \cdot d{\bf r}, \,\, \int_{C_3} {\bf F} \cdot d{\bf r}\right)}$ from smallest to largest.  Explain your reasoning.
%
%\hspace{11.0cm} \includegraphics[scale = 0.4]{../../../Figures/Test_Fig15.png}
%
%\vspace{0.4cm}
%
%(b) Suppose \({\bf F}\) is conservative and 
%\(\int_{(0,0)}^{(0,1)}{\bf F} \cdot d{\bf r} = 2.4\), \,\,
%\(\int_{(0,1)}^{(1,1)}{\bf F} \cdot d{\bf r} = -1.2\) \,\, and
%\(\int_{(0,0)}^{(1,0)}{\bf F} \cdot d{\bf r} = 7.5\). \,\,Then
%
%$\D{\int_{(1,0)}^{(1,1)}{\bf F} \cdot d{\bf r} = }$
%
%\vspace{0.4cm}
%
%(c) Let ${\bf F} = [-y,x]$ and $C$ be a simple closed curve in the $xy-$plane. 
%
%Show that $\D{\oint \limits_{C} \text{\bf F} \cdot d\text{\bf r} = 2\cdot \text{Area}(R)}$ \quad where $R$ is the region bounded by $C$.
%
%\vspace{0.2cm}
%
%\thispagestyle{empty}
%\newpage
%\thispagestyle{empty}

%{\bf Problem 6:} Let $\text{\bf F} = [x+y,2xy]$ and $C$ be the boundary of the region enclosed by the curves $y=x^2$, $y=x+2$, and $x=0$ in the first quadrant. \\
%(a) \UL{Setup} the integral you would use to compute \,\, $\D{\oint \limits_{C} \text{\bf F}\cdot d\text{\bf r}}$ \,\, directly \,\,\,({\em don't leave integrand as a dot product})
%
%(b) Compute \,\, $\D{\oint \limits_{C} \text{\bf F}\cdot d\text{\bf r}}$ \,\, using Green's theorem.
%
%\thispagestyle{empty}
%\newpage
%\thispagestyle{empty}
%
%Let $f(x,y,z)$ be a scalar-valued function.  Using the concepts from this class, explain why curl$(\D \nabla f) = {\bf 0}$ no matter what $f(x,y,z)$ is.  If you cannot explain it, then prove it mathematically.
%
%%\includegraphics[scale = 0.4]{../../../Figures/Test1_Fig4.png}
%
%%\includegraphics[scale = 0.4]{../../../Figures/Test1_Fig7.png}
%
%\includegraphics[scale = 0.4]{../../../Figures/Test1_Fig6.png}

\end{document}